\textbf{III -- LA STRUTTURA DEI PEDONI}

\textbf{1. L'indebolimento delle case e il ``buco'' (``foro'',
``hole'')}

Nel commento alla mossa 12 del Nero nella partita Spassky-Fischer, il
campione americano diceva di avere scartato una certa mossa perché
avrebbe condotto a una posizione piena di «buchi». Che cosa intendeva
dire Fischer di preciso? È quel che cercheremo di chiarire in questa
lezione.

Nella posizione iniziale i giocatori controllano ognuno metà scacchiera
e nessun pezzo può installarsi in modo duraturo oltre la terza traversa
perché basta una semplice mossa di pedone per scacciarlo. Se togliessimo
tutti i pedoni dalla scacchiera i pezzi leggeri potrebbero occupare,
senza temere pe¢dite di tempo, anche posizioni avanzate; un eventuale
attacco non li costringerebbe a retrocedere perché potrebbero essere
difesi dagli altri pezzi. Si può dunque affermare che il controllo che
un pedone esercita su una casa è più forte di quello esercitato da un
pezzo. Ne consegue che la funzione delle catene di pedoni è quella di
guadagnare spazio per i pezzi che così possono svilupparsi indisturbati
dietro le linee.

Ovviamente non è tutto oro quel che riluce. Le spinte di pedoni
guadagnano territorio ma a scapito della sicurezza delle retrovie; le
case che si lasciano alle spalle si indeboliscono in quanto non più
difendibili in modo forte (con un pedone). ``Ogni spinta di pedone crea
una debolezza'' sentenziava Steinitz. Il vantaggio di spazio in genere
compensa l'indebolimento delle retrovie ma occorre evitare certe
configurazioni di pedoni che rendono questo indebolimento
particolarmente grave, tanto da diventare, in certi casi, il motivo
dominante della partita. La struttura dei pedoni fa da scudo non solo ai
pezzi ma anche alle case arretrate, le ripara dalle intrusioni
avversarie. Se però nella struttura si producono squarci le case interne
possono indebolirsi a tal punto da diventare buchi.

Ora dovrebbe essere chiaro che cosa si intende per buco: un buco nello
schieramento, poniamo del Bianco, è una casa, nella sua zona di
interesse, non controllabile dai pedoni bianchi; ciò significa che i
pedoni delle colonne adiacenti o sono assenti o sono troppo avanzati.
Chiamiamo zona di interesse di un colore l'insieme delle case occupate,
adiacenti o retrostanti, a pedoni di quel colore. È ovvio che la
presenza del buco sarà ancora più grave se esso è controllato da pedoni
avversari, ma non è detto che un buco renda una casa necessariamente
debole. L'uso della parola è comunque oscillante nei vari autori; da un
lato può estendersi alle case deboli in genere, dall'altro restringersi
al caso sopra ricordato del controllo da parte di un pedone avversario.

Vediamo due esempi:

\bookdiagramplaceholder{assets/diagrams/image21.png}{21}

Nello schieramento nero c'è un buco in d5.

\bookdiagramplaceholder{assets/diagrams/image22.png}{22}

Osserviamo la posizione e facciamo una piccola ricerca di ``buchi''.

Nello schieramento nero se ne contano diversi; in particolare in a6, c6,
f5, f6 e h5. È ovvio che non tutte le case elencate hanno la stessa
importanza strategica. La casa a6 è trascurabile perché decentrata e lo
stesso dicasi per la casa c6 se il Nero ha effettuato l'arrocco corto.
In caso di arrocco lungo sarebbe appetibile perché a ridosso del Re, ma
comunque di difficile sfruttamento se protetta dall'Alfiere in
fianchetto (in b7).

Passiamo a f6: anche senza conoscere la disposizione dei pezzi il piano
di piazzarcene uno in pianta stabile sembra pretenzioso: un Cavallo
potrebbe arrivarci solo attraverso h5, casa che può essere facilmente
guardata con un Cavallo nero posto in g7 o in f6; una Torre, invece,
potrebbe essere scacciata con facilità, per esempio da un Cavallo in d7.

Il buco classico, quello che attira l'attenzione del giocatore esperto,
si trova in f5. Qui si ha una situazione ideale. La casa f5 è per il
Bianco una casa ``forte'', controllata com'è da due pedoni e, per di
più, è facilmente accessibile ai pezzi. Può esserci installato un
Cavallo facendolo passare da g3, casa questa dalla quale i pezzi bianchi
non possono essere in alcun modo intercettati. L'avamposto in f5 può
essere facilmente rafforzato da un secondo Cavallo in g3 e dalle Torri,
eventualmente raddoppiate sulla colonna f.

L'unico foro nello schieramento bianco si trova in h4 ma, poiché
laterale, non è strategicamente di gran valore.

Le conseguenze derivanti anche da un solo buco possono portare alla
sconfitta.

Vediamo che cosa successe nella seguente partita:

Janowsky--Kupchik, L'Avana 1913.

\textbf{1.d4 d5 2.c4 e6 3.¤c3 ¤f6} per evitare continuazioni legate allo
sviluppo del Cavallo in e2 alcuni preferiscono far precedere questa
mossa da 3..¥e7; impedito 3.¥g5, il Bianco non ha modo di ritardare lo
sviluppo del Cavallo di Re.

\textbf{4.¥g5 ¥e7 5.e3 ¤bd7 6.¥d3 £xc4} nella partita di Donna spesso si
assiste a una lotta per il tempo legata alla presa in c4. In molte
varianti il Bianco aspetta che il Nero abbia preso in c4 per uscire
d'Alfiere e il Nero aspetta che il Bianco abbia mosso l'Alfiere per
prendere in c4.

\textbf{7.¥xc4 ¤b6 8. ¥d3 ¤fd5 9.¥xe7 £xe7 10.¤f3 0-0 11.0-0 ¥d7 12.¦c1
c6 13.¤e4}

\bookdiagramplaceholder{assets/diagrams/image23.png}{23}

Si è raggiunta una posizione in cui il Bianco è uscito decisamente
meglio dall'apertura. Ha completato in modo sano lo sviluppo e ora può
intraprendere le operazioni di attacco sull'arrocco nero, così decide di
portare un altro pezzo nella zona delle operazioni. Il Nero ha ancora da
completare lo sviluppo: l'¥d7 è accecato dai propri pedoni, il ¤b6 è per
il momento fuori gioco. Intimorito dalle operazioni che il Bianco
minaccia di intraprendere sul lato di Re, il Nero decide di chiudere la
pericolosa diagonale b1-h7 nel modo più drastico:

\textbf{13... f5?} e commette una mossa suicida. In e5 si è infatti
creato un buco che il Bianco non tarderà ad occupare per farne il perno
del proprio attacco.

\textbf{14.¤c5 ¥e8} per portarlo in h5 e dare respiro al povero Alfiere.

\textbf{15.¤e5 ¦b8} così facendo il bianco evita 16... ¥h5 e va nella
casa debole. L'idea del Nero è di giocare ¤d7 per barattare il Cavallo
con quello in e5. Naturalmente il Bianco si guarderebbe bene
dall'acconsentire il cambio; egli è pronto a eliminare un eventuale
Cavallo in d7 col proprio c5. In tal caso otterrebbe anche di riportare
l'Alfiere campo chiaro del Nero nella disgraziata casa d7.

\textbf{16.¦e1 ¦f6} il Nero cerca disperatamente un controgioco sul lato
di Re con ¦h6 e £h4.

\textbf{17.£f3 ¦h6 18.£g3} minaccia 19.¤f7 che attacca
contemporaneamente la ¦h6 con ¤f7 e la ¦b8 con la Donna.

\textbf{18... ¦c8} ammettendo che la quindicesima mossa è stata una
pe¢dita di tempo.

\textbf{19.f3 ¦c7} il Bianco prepara la spinta in e4.

\textbf{20.a3 ¢h8 21.h3 g5 22.e4} ora che tutti i pezzi sono ben
piazzati il Bianco apre il gioco.

\textbf{22... f4} le spinte di pedone hanno creato debolezze nello
schieramento bianco e anche il Nero ha conquistato una casa in
territorio nemico. Questa mossa prepara la discesa di un pezzo.

\textbf{23.£f2 ¤e3} il Cavallo si installa nel buco e3, in posizione
piuttosto fastidiosa.

\textbf{24.¦xe3!} Janowsky non esita a sacrificare la qualità, in cambio
di un solo pedone, pur di eliminare l'avamposto del Nero. Il sacrificio
è più che giustificato perché elimina l'unico pezzo attivo del Nero. Ora
il Bianco ha le mani libere per portare avanti l'attacco. Per qualità si
intende la differenza che intercorre tra il valore della Torre e quella
di un Alfiere o di un Cavallo.

\textbf{24... fxe3 25.£xe3 ¤c8 26.¤g4}

\bookdiagramplaceholder{assets/diagrams/image24.png}{24}

L'idea di questa mossa è di portare il Cavallo in f6. Attaccando la
Torre il Bianco guadagna il tempo necessario per spingere in e5 e
controllare la casa f6. Janowsky rinuncia al buco in e5 per quello, più
avanzato, in f6.

\textbf{26... ¦g6 27.e5 ¦g7 28.¥c4 ¥f7 29.¤f6 ¤b6} questo Cavallo
ballonzola da una casa all'altra senza posa né costrutto.

\textbf{30.cce4} rafforza la presenza del Cavallo in f6.

\textbf{30... h6 31.h4 ¤d5 32.£d2 ¦g6} il Nero è disposto a restituire
la qualità pur di scacciare l'indesiderato intruso.

\textbf{33.hxg5! £f8} la presa 33... hxg5? è sbagliata per via di 34.¢f2
seguita da ¦h1, con facile vittoria.

\textbf{34.f4 ¤e7 35.g4 hxg5 36.fxg5} il Nero abbandona.

Il Bianco minaccia matto in due con 37.£h2+ ¢g7 38.£h7 matto. A 36...
¥g8, che difende in h7, segue 37.£h2+ ¢g7 38.¥xe6~.

Si rifletta come le case e5 prima e f6 dopo siano state per il Nero
delle vere spine nel fianco.

\textbf{2. I pedoni deboli}

Un pedone si dice debole quando non può essere difeso da altri pedoni.
Come è facile immaginare ci sono vari casi di pedoni deboli e vari gradi
di debolezza.

\textbf{a) Il pedone isolato}

Se, nelle colonne adiacenti a quella in cui si trova un pedone, non ce
ne sono altri dello stesso colore (perché catturati o cambiati), quel
pedone si dice ``isolato''.

\bookdiagramplaceholder{assets/diagrams/image25.png}{25}

Il Pedone in d5 è isolato.

È fuor di dubbio che il pedone isolato sia più debole degli altri ma,
anche se bisogna averne una cura particolare, non sempre è svantaggioso
averne uno.

Può infatti intervenire una compensazione di altra natura, in genere
legata a una maggiore attività dei pezzi leggeri. Per esempio, nella
posizione 25, gli Alfieri neri non trovano ostacoli nei propri pedoni
mentre uno bianco piazzato, poniamo, in c1 avrebbe meno libertà di
azione per via del Pe3. Non stupisca quindi che alcuni sistemi
d'apertura prevedano la formazione di un pedone isolato. Però
attenzione: quando ciò accade il pedone è anche mobile (cioè può essere
spinto ancora). L'ulteriore spinta ne assicura il cambio prima che
diventi debole davvero e può essere usata come ariete contro la catena
di pedoni avversaria. Se, invece, un pedone isolato è bloccato allora la
debolezza è maggiore e può farsi sentire anche in modo pesante durante
la partita.

La debolezza è più grave ancora se, oltre ad essere bloccato, il pedone
isolato si trova su colonna semiaperta (intendendo con questo termine
una colonna sgombra di altri pedoni) perché può essere attaccato anche
con le Torri, eventualmente raddoppiate proprio su quella colonna. Un
simile pedone può facilmente diventare un obiettivo strategico di
primaria importanza.

La formazione di un pedone isolato muta l'indirizzo strategico della
partita: \emph{``Il pedone isolato getta un'ombra sinistra su tutta la
scacchiera''}, disse Tartakower.

Si capisce allora quale strategia adottare appena un pedone isolato si
manifesta nella posizione avversaria. Per prima cosa lo si blocca,
concentrando le forze nella casa successiva a quella in cui il pedone è
posto. Consolidata la conquista della casa di blocco, la si occupa con
un pezzo. Bisogna fare in modo di riprendere, in caso di cambio, con un
altro pezzo e non di pedone. Un'eventuale ripresa col pedone chiuderebbe
la colonna e renderebbe il pedone isolato meno debole. Infine si
mobilitano le forze per farlo cadere.

Vediamo, in due esempi tratti dalle teoria delle aperture, come si opera
per bloccare un pedone isolato.

\emph{Difesa Siciliana:} \textbf{1.e4 c5 2.c3} questa mossa conosce da
qualche anno una certa popolarità, soprattutto perché evita di studiare
le numerose e complicate varianti della Siciliana aperta (2.¤f3 3.d4).

L'obiettivo è di guadagnare spazio al centro, utile per imbastire un
eventuale attacco sull'ala di Re dove è presumibile il Nero arroccherà,
ma ha lo svantaggio di creare un pedone isolato nel proprio
schieramento.

\textbf{2... d5} l'alternativa più comune, 2... ¤f6, evita di portare
subito in gioco la Donna che potrà essere attaccata con pe¢dita di tempo
ma non provoca la formazione di un pedone isolato.

\textbf{3.exd5 £xd5 4.d4 e6} il Nero può rita¢dare il cambio in d4 in
quanto il Bianco non può giocare £xc5 per via di £xd1+.

\textbf{5.¤f3 ¤f6 6.¥d3 ¤xd4 7.¤xd4} il Bianco ha un pedone isolato in
d4. Le prossime mosse del Nero sono finalizzate a consolidare il
controllo della casa d5.

\textbf{7... ¥e7 8.¤c3 £d8 9.0-0 0-0 10.£e2 ¤c6 11.¦d1 ¤b4!} attacca
l'¥d3 e guadagna il tempo necessario per occupare la casa d5.

\textbf{12.¥c4 b6}

\bookdiagramplaceholder{assets/diagrams/image26.png}{26}

Con l'idea di piazzare l'Alfiere in fianchetto e consolidare
ulteriormente la casa d5. Le mosse successive del Nero saranno dunque
¥b7 e ¤d5, per bloccare in modo definitivo il Pd4 e farne in seguito
oggetto di attacco.

\emph{Difesa Francese:} \textbf{1.e4 e6 2.d4 d5 3.¤d2 c5 4.exd5 exd5
5.¤f3 ¤c6 6.¥b5 ¥d6 7.0-0 ¤ge7 8.£xc5 ¥xc5} in d5 si è creato un pedone
isolato. Dopo quanto si è detto lo scopo delle successive mosse del
Bianco dovrebbe risultare chiaro.

\textbf{9.¤b3 ¥d6 10.c3 ¥g4 11.¤bd4}

\bookdiagramplaceholder{assets/diagrams/image27.png}{27}

Si noti come in entrambi gli esempi chi blocca abbia superprotetto il
pezzo bloccante in modo da evitare di riprendere con un pedone nel caso
venga cambiato.

\textbf{b) Il pedone arretrato}

Un caso molto simile al pedone isolato, e per il quale valgono molte
delle considerazioni già fatte, è quello del pedone arretrato. Si dice
arretrato un pedone che non può essere difeso da altri pedoni e che si
trova bloccato; il blocco è anche qui più efficace se la casa davanti al
pedone è guardata da almeno un pedone avversario. Anche in questo caso
se il pedone si trova su colonna semiaperta è particolarmente debole.

Il pedone d6 è debole perché arretrato.

\bookdiagramplaceholder{assets/diagrams/image28.png}{28}

Vediamo ora che cosa succede in una variante della difesa Siciliana
venuta in una certa moda nell'ultimo ventennio per merito del GM russo
Evgeni Sveshnikov, da cui prende il nome:

\textbf{1.e4 c5 2.¤f3 ¤c6 3.d4 ¤xd4 4.¤xd4 ¤f6 5.¤c3 e5}

\bookdiagramplaceholder{assets/diagrams/image29.png}{29}

Prima degli anni Settanta nessuno giocava così e in tutta la storia
degli scacchi potevano contarsi pochissime partite con questo o¢dine di
mosse. Il motivo risiede nelle debolezze che la mossa e7-e5 provoca
nella struttura dei pedoni neri. Poiché il pedone sulla colonna \emph{c}
non c'è più, la spinta del pedone \emph{e} dà origine inevitabilmente a
debolezze. In caso di spinta di un passo si indebolisce il Pd6 e, se
spinto di due passi, si forma anche un buco in d5. Queste
considerazioni, peraltro giuste, portavano a scartare a priori una mossa
che creava tanti problemi al Nero. Si pensi che nel \emph{``Manuale
teorico pratico delle aperture''} di Giorgio Porreca di 772 pagine
(Milano, Mursia), edito per la prima volta nel 1971 ma ancora in
commercio in una ennesima ristampa, a tale mossa, oggi entrata nel
repertorio anche di grandi maestri di primo piano, è dedicata solo una
piccola nota in cui, peraltro, si dà una continuazione sbagliata.
Sveshnikov fu il primo a rendersi conto che il Nero otteneva vantaggi di
altra natura (grande mobilità dei pezzi leggeri, possibilità di
contrattacco, rimedio drastico contro la spinta di rottura e4-e5) ma la
sua opera fu vista con sufficienza a lungo. Io sono stato uno dei primi
ad adottarla in Italia.

Nel 1977 Sveshnikov venne a giocare un torneo in Italia (Marina Romea).
Per caso alloggiavamo nello stesso albergo e la sera dava spettacolo
vincendo una ``lampo'' dietro l'altra a un maestro russo. Io (allora
semplice Terza Nazionale) rimasi colpito dalla sua variante e decisi di
farne la mia difesa. Vantaggi e svantaggi di questa difficile apertura
sono ben mostrati nella partita che giocai nel Magistrale di Bagni di
Lucca del 1984, contro il maestro americano Stuart Wagman.

Wagman--Leoncini, Bagni di Lucca 1984.

\textbf{1.e4 c5 2.¤f3 ¤c6 3.d4 ¤xd4 4.¤xd4 ¤f6 5.¤c3 e5} oltre che col
nome di Sveshnikov in certi manuali questa variante chiamata anche
Lasker e, di rado, Pilnik o Cebljansky.

\textbf{6. ¤d5 d6} la posizione è caratterizzata da un buco in d5 e un
pedone debole in d6.

\textbf{7.¥g5 a6 8.¤a3 b5 9.¤d5} il Bianco occupa la casa debole.

\textbf{9... ¥e7 10.¥xf6 ¥xf6 11.c3 0-0 12.¤c2} un secondo Cavallo si
prepara a dar man forte al primo via c2 ed e3.

\textbf{12... ¥g5} a 12... ¦b8 il Bianco replica con 13.h4! che prepara
l'attacco sull'arrocco e toglie la casa g5 all'Alfiere.

\textbf{13.a4} l'idea di questa mossa è di distogliere il Pb5 dal
controllo della casa c4 dove il Bianco intende piazzare un Alfiere. Un
Alfiere in c4 fortifica ulteriormente la casa d5 e fa pressione
minacciosamente sulla diagonale a2-g8.

\textbf{13... bxa4 14.¦xa4 a5 15.¥c4 ¦b8} pone sotto pressione il pedone
debole b2.

\textbf{16.¦a2 ¢h8} se il Nero vuole riconquistare la casa d5 deve
minare innanzitutto il Pe4. Questa mossa, che toglie il Re dalla
diagonale a2-g8, prepara la spinta in f7-f5.

\textbf{17.cce3 ¥xe3?!} una mossa di dubbio valore. Il Nero elimina un
controllore di d5 ma il cambio dei pezzi leggeri diminuisce le
possibilità di controgioco.

\textbf{18.¤xe3 ¤e7 19.0-0 f5?} il piano è giusto ma eseguito troppo in
fretta. Il Nero punta a indebolire il Pe4 e ad attaccare sul lato di Re,
ma occorreva far precedere questa mossa da g7-g6 per rispondere a exf5
con gxf5 ed evitare un ulteriore cambio di pezzi leggeri.

\textbf{20.exf5 ¤xf5 21.¤xf5 ¥xf5} questi cambi hanno spuntato le
possibilità di controgioco del Nero e solo un errore del Bianco può
salvarlo dalla sconfitta.

\textbf{22.b3} consolidata la posizione il Bianco è pronto ad iniziare
l'attacco diretto sui pedoni deboli in a5 e d6.

\textbf{22... £d6 23.£d5!} molto ben giocato. Questa mossa mette sotto
osservazione entrambi i pedoni da attaccare e prepara il ra£doppio sulla
colonna con le Torri.

\textbf{23... ¦a8 24.¦d1 ¦fd8} il Nero è costretto a giocare in modo
passivo.

\textbf{25. ¦a1 ¦f8} con l'idea di rispondere con ¥e4 alla presa in a5.
Ormai il Nero è costretto ad arrampicarsi sugli specchi per tenere in
piedi la posizione.

\textbf{26.¦xa5?}

\bookdiagramplaceholder{assets/diagrams/image30.png}{30}

Ecco l'errore! Il Bianco, che fino a questo punto aveva giocato in modo
impeccabile, ora pecca di precipitazione. Wagman aveva visto la replica
del Nero (28... ¥e4) ma aveva preparato una sorprendente mossa di Donna,
purtroppo per lui non sufficiente a contenere il controgioco. Si doveva
far precedere questa presa da 26.£d5! Ora il Nero prende in mano le
redini della partita.

\textbf{26... ¥e4!} minaccia matto in due (27... £xf2+ 28.¢h1 £xg2
matto); se 27.£xe4 ¦xa5.

\textbf{27. £f7!} Wagman aveva riposto le sue speranze in questa mossa
spettacolare; in effetti £f7 salva il materiale ma espone la Donna a
ripetuti attacchi. Con i tempi guadagnati il Nero imbastirà un attacco
micidiale.

\textbf{27... ¦ac8 28.¦b5} intermedia necessaria per scacciare la Donna
dalla diagonale a7-g1.

\textbf{28... £c6} al solo costo di un pedone il Nero ha conquistato
l'iniziativa e ricolloca i pezzi sulla scacchiera imbastendo minacce su
minacce.

\textbf{29.£e7 ¦ce8} se 29... ¥xg2? 30.¦a7! ¦g8 31.¥xg8 ¦xg8 3 2.¦b8! e
matto alla successiva.

\textbf{30.£g5 ¦f6} la minaccia 31... ¦g6 costringe il Bianco a
indebolire l'arrocco.

\bookdiagramplaceholder{assets/diagrams/image31.png}{31}

\textbf{31.g3
d5 32.¥f1}

\textbf{32... ¦xf2!! 33.¢xf2 ¦f8+ 34.¢g1} a 34. ¢e3 (o 34.¢e1) segue
34... £xc3+.

\textbf{34... £xc3 35.}~il Bianco abbandona. Non c'è modo di parare
tutte e due le minacce 34... £d4+ e 34... £f3 con matto alla seguente.

Da tutto ciò vediamo di trarne qualche insegnamento. Avere ``buchi'' e
pedoni deboli nel proprio schieramento è uno svantaggio che, se non
compensato da altri fattori, conduce alla sconfitta. Quindi occorre
evitare di crearsi debolezze se non si è più che sicuri di ottenere
vantaggi di altra natura.

\textbf{c) Pedoni doppiati (o impedonati)}

Con questo termine si indicano i pedoni del solito colore posti sulla
stessa colonna. La doppiatura dei pedoni è uno svantaggio minimo, di
certo meno grave di un pedone arretrato, anche perché di solito
compensato dall'apertura delle linee.

Gli svantaggi sono intuibili, basti per tutte questa semplice posizione:

\bookdiagramplaceholder{assets/diagrams/image32.png}{32}

Come si vede il Bianco non ha la possibilità di sfruttare il pedone in
più.

La qualità dello svantaggio di una doppiatura è di diversa portata nelle
diverse fasi della partita e tende ad aumentare a mano a mano che ci si
avvicina al finale. In fase di sviluppo può trovare sufficiente compenso
dall'apertura delle linee; nel centro partita può essere utile avere una
colonna aperta in più ma occorre fare attenzione a che non si trasformi,
con un semplice baratto, in pedone debole; in finale la doppiatura può
equivalere, come nell'esempio, ad avere un pedone in meno.

I pedoni doppiati irrigidiscono le catene di pedoni e facilitano lo
sfondamento avversario; ecco un esempio classico tratto dalla teoria dei
finali:

\bookdiagramplaceholder{assets/diagrams/image33.png}{33}

Il Bianco vince giocando 1.h5 (1.g5?? fxg5 2.fxg5 h5 e vince il Nero).
seguito da 2.g5.

L'idea di provocare un'impedonatura, e di giungere rapidamente in finale
per sfruttarla a pieno, si ha nella variante di Cambio della Partita
Spagnola:

\textbf{1.e4 e5 2.¤f3 ¤c6 3.¥b5 a6 4.¥xc6 £xc6 5.d4 exd4 6.£xd4 £xd4
7.¤xd4}

\bookdiagramplaceholder{assets/diagrams/image34.png}{34}

Se ora nella scacchiera lasciassimo solo i pedoni il Bianco vincerebbe
facilmente grazie al quattro contro tre sul lato di Re mentre il pedone
in più nero sul lato di Donna, a causa dell'impedonatura, equivale a un
tre contro tre. Al finale di pedoni occorre naturalmente arrivarci; in
questa particolare variante il Nero ha compenso sufficiente grazie alla
coppia degli Alfieri.

\emph{Difesa Francese:} \textbf{1.e4 e6 2.d4 d5 3.¤c3 ¥b4 4.e5 c5 5.a3
¥xc3 6.bxc3}

\bookdiagramplaceholder{assets/diagrams/image35.png}{35}

Anche qui l'impedonatura ha, come corrispettivo, la coppia degli
Alfieri. Notare che il pedone in c2 del Bianco può diventare debole e
arretrato dopo il cambio in d4.

\emph{Difesa Caro Kann:} \textbf{1.e4 c6 2.d4 d5 3.¤c3 £xe4 4.¤xe4 ¤f6
5.¤xf6 exf6}

\bookdiagramplaceholder{assets/diagrams/image36.png}{36}

Il Nero ha concesso al Bianco un finale superiore grazie al quattro
contro tre sul lato di Donna. Per compenso ha ottenuto uno sviluppo
scorrevole e un rafforzamento del lato di Re dove intende arroccare. A
differenza delle varianti precedenti il Bianco, però, non si è privato
della coppia degli Alfieri.

I manuali sulle aperture mettono in evidenza che l'idea della difesa
Nimzo Indiana è una sorta di lotta per il controllo delle case centrali.

\textbf{1.d4 ¤f6} questo Cavallo impedisce di piazzare un secondo pedone
al centro e gua¢dare anche la casa d5.

\textbf{2.c4} il Bianco attua un'espansione sul lato di Donna e gua¢da
lo stesso la casa d5.

\textbf{2... e6} stabilisce un secondo controllo su d5 contro uno solo
del Bianco.

\textbf{3.¤c3} pareggia il conto su d5 e minaccia e2-e4.

\textbf{3... ¥b4} inchiodando il ¤c3 annulla gli effetti dell'ultima
mossa bianca.

\textbf{4.a3} la variante Sämisch.

\textbf{4... ¥xc3+ 5.bxc3} nella struttura di pedoni del Bianco è
comparsa una doppiatura di pedoni che può diventare obiettivo del Nero.
Il Nero dovrà prima bloccarlo con la spinta c7-c5 e poi attaccarlo con
¤c6-a5 e ¥c8-a6 (dopo avere ovviamente spinto in b6. Un buon esempio di
questa strategia si ebbe nella partita Geller-Smyslov, giocata
all'Interzonale di Amsterdam nel 1956. Eccone il seguito:

\textbf{5... c5 6.e3 b6 7.¤e2 ¤c6 8.¤g3 0-0 9.¥d3 ¥a6 10.e4 ¤e8 11.¥e3
¤a5 12.£e2 ¦c8 13.d5 £h4 14.0-0 ¤d6} e non ci sono più difese per il
Pc4.

Da questo semplice esempio non si commetta l'errore di credere che la
variante Sämisch sia demolita. Per il pedone doppiato il Nero ha ceduto
la coppia degli Alfieri e si è indebolito il lato di Re. Tanto per dare
un'idea delle possibilità bianche in questa variante si veda la partita
che segue, in cui il lato di Re nero viene letteralmente distrutto.

Lilienthal--Najdorf, Stoccolma 1948, \textbf{(Difesa Nimzo Indiana)}

\textbf{1.d4 ¤f6 2.c4 e6 3.¤c3} per evitare la Nimzo Indiana si può
giocare 3.¤f3 b6 (3... ¥b4+ 4.¥d2) 4.a3 ¥b7 5.¤c3 d5 6.¤xd5 ¤xd5 7.e3
ecc.

\textbf{3... ¥b4 4.a3} le alternative principali, e più comuni di questo
tratto, sono 4.e3 (variante Rubinstein) e 4.£c2 (variante Capablanca).

\textbf{4... ¥xc3+ 5.bxc3 c5 6.e3 ¤c6 7.¥d3 b6} l'occupazione del centro
può essere contrastata con 7... e5 8.¤e2 b6 9.0-0 0-0 10.e4 h6 11.d5 ¤e7
12.f3 ¤g6 13.¥e3 ¥d7 14.£d2 ¦b8 15.a4 b6 16.¤g3 con gioco all'incirca
equivalente, come nella partita Balashov-Lerner, campionato dell'URSS
1985.

\textbf{8.¤e2 0-0 9.e4 ¤e8 10.0-0 d6} è migliore 10... ¥a6 11.f4 f5
12.¤g3 ¤e7 13.exf5 ¤xf5 14.¤xf5 exf5 (Martin-Short, Londra 1985).

\textbf{11.e5!} in considerazione che non è possibile 11.¤xe5?? per
12.¥xh7+ ¢xh7 13.£xd8.

\textbf{11... £xe5 12.£xe5 ¥b7} il vantaggio di spazio del Bianco, che è
ormai pronto a dar l'assalto all'arrocco, non permette al Nero di
sfruttare le debolezze nella struttura di pedoni bianca; anzi, in questa
fase l'impedonatura svolge un ruolo positivo perché permette il
controllo delle case centrali d5 e e5.

\textbf{13.¥f4 f5} un tentativo disperato, ma dopo 13... £c7 il Nero non
ha alcun controgioco da contrapporre all'iniziativa dell'avversario ad
est.

\textbf{14.exf6 e5?} era meglio 14... £xf6 ma al campione a¢gentino era
sfuggita la forza del sacrificio d'Alfiere.

\textbf{15.fxg7 ¦xf4 16.¤xf4 exf4}

\bookdiagramplaceholder{assets/diagrams/image37.png}{37}

\textbf{17.¥xh7+!} completa la demolizione dell'arrocco nero. Ora il Re
cade in balìa dei pezzi bianchi.

\textbf{17... ¢xh7 18.£h5+ ¢xg7 19.¦ad1 £f6 20.¦d7+ ¢f8 21.¦xb7 ¤d8
22.¦d7 ¤f7 23.£d5 ¦b8 24.¦e1 f3} che altro? A una maggiore resistenza
avrebbe portato 24... ¤e5 25.£xa8 ¤xd7 26.¦e1 £g6 per poi spingere in f3
ma al Bianco è sufficiente collocare la Donna sulla grande diagonale
bianca per evitare il matto.

\textbf{25.¦e3 abbandona}

\begin{enumerate}
\def\labelenumi{\alph{enumi})}
\setcounter{enumi}{3}
\item
  \textbf{Pedoni sospesi}
\end{enumerate}

\bookdiagramplaceholder{assets/diagrams/image38.png}{38}

I pedoni c5 ed d5 sono deboli perché su colonne aperte. La spinta di uno
di essi provoca il cambio e la debolezza del pedone rimanente. Però,
finché sulla scacchiera sono presenti molti pezzi, i pedoni sospesi
centrali possono rappresentare un punto di forza; la maggiore libertà di
azione di cui godono i pezzi posti dietro ai pedoni passati possono
creare le condizioni per una veloce avanzata verso la promozione di uno
di essi, come accade in questa partita.

Sokolsky--Botvinnik, Leningrado 1938 \textbf{(Difesa Grunfeld)}

\textbf{1.c4 ¤f6 2.¤c3 d5 3.d4 g6 4.¤f3 ¥g7 5.e3 0-0 6.¥e2 e6 7.0-0 b6
8.¤xd5} poiché il Nero si appresta a sviluppare l'Alfiere in fianchetto,
Sokolsky si affretta a chiudere la diagonale a8-h1, ma secondo Botvinnik
era meglio mantenere la tensione al centro con 8.£d3.

\textbf{8... exd5 9.b3 ¥b7 10.¥b2 ¤bd7} per controllare le case e5 e c5.

\textbf{11.£c2 a6 12.¦ac1 ¦c8 13.¦fd1 £e7 14.£d1 ¦fd8 15.¥f1 c5 16.£xc5}
commenta Botvinnik: \emph{``Un altro strafalcione posizionale. La
debolezza dei pedoni neri in d5 e c5 non può essere sfruttata, quando è
presente un gran numero di pezzi leggeri e i pedoni sono attaccati dalla
prima traversa. Il Bianco cambia invece il §d4, cedendo l'ultimo
caposaldo centrale; ciò ridà vita all'¥b7''.}

\textbf{16... bxc5}

\bookdiagramplaceholder{assets/diagrams/image39.png}{39}

\textbf{17.¤e2 ¥h6} minaccia la spinta in d4.

\textbf{18.¥a3} per rispondere a 18... d4 con 19.¦xd4.

\textbf{18... ¤g4 19.£d3 C£e5 20.¤xe5 £xe5 21.¤g3 £f6 22.ch1 d4 23.£e2
¤e5 24.exd4 ¤xd4}

\bookdiagramplaceholder{assets/diagrams/image40.png}{40}

Uno dei due pedoni si è trasformato in pedone passato (o libero),
intendendo con questo termine un pedone la cui marcia non è ostacolata
da pedoni avversari. Se il Bianco avesse modo di bloccarlo, ridisposte
le forze in campo, potrebbe in seguito attaccarlo.

Ma così non è. A causa della forte attività dei pezzi neri il Bianco non
ha modo di bloccarlo ed esso in pochi tratti marcerà trionfalmente a
Donna.

\textbf{25.¦xc8 ¥xc8} la ¦d8 non va distolta dal Pedone libero.

\textbf{26.¦e1 d3! 27.£d1} se 27.£xe5? £xe5 29.¦xe5 d2 e vince.

\textbf{27... ¥g4 28.£a1 d2 29.¦xe5 d1=£ 30.¦e8+ ¦xe8 31.£xf6} finite le
prese il Nero è rimasto con la qualità in più, di certo non compensata
da un unico pedone. Il resto è questione di tecnica.

\textbf{31... ¥e2 32.¤g3 ¥g7 33.£c6 ¥b5 34.£c1 £xc1 35.¥xc1 ¦e1 36.¥e3
¦a1 37.a4 ¥d3 38.f4 ¦b1 39.¢f2 ¥xf1 40.¤xf1 ¦xb3 41.abbandona}

\textbf{3. L'attacco di minoranza}

I cambi possono con facilità creare una asimmetria nella configurazione
pedonale dei giocatori; può crearsi una maggioranza di pedoni su un lato
cui fa da contrappunto quella avversaria sul lato opposto.

La pratica ha dimostrato che se la maggioranza di pedoni è poco mobile
può essere vantaggiosamente attaccata dalla minoranza di pedoni. Quello
che a prima vista può sembrare un paradosso si capisce se si è
familiarizzato con la teoria dei pedoni deboli.

In effetti lo scopo dell'attacco di minoranza è proprio di creare un
pedone debole (sia esso isolato o arretrato) nello schieramento
avversario. Il motivo per cui questo è possibile non è difficile da
capire: cambiando i pedoni ne rimarrà uno, su colonna semiaperta, che
inevitabilmente sarà debole.

\bookdiagramplaceholder{assets/diagrams/image41.png}{41}

Questa è la posizione classica in cui il Bianco può operare un attacco
di minoranza. Sul lato di Donna il Nero ha una maggioranza di quattro
pedoni contro solo tre bianchi, ma la loro mobilità è ridotta. Infatti,
la spinta in c5 isolerebbe il Pd5 dopo d4xc5; la spinta in b6 renderebbe
più debole quello in c6. In queste condizioni un eventuale assalto del
Bianco è giustificato. L'attacco avviene spingendo a fondo i pedoni
``a'' e ``b''. Giunto in b5 il pedone minaccia la presa in c6 perché
dopo b7xc6 nella struttura del Nero si forma un pedone arretrato (c6).
D'altra parte, se il Nero gioca d'anticipo c6xb5, dopo a4xb5 il Bianco
può puntare i suoi cannoni sul pedone ``a''. Infine, come abbiamo già
detto, non serve nemmeno la spinta c6-c5 perché dopo d4xc5 il pedone d5
rimane isolato.

Prima di esaminare l'attacco di minoranza nella pratica, vediamo al
solito, alcune aperture che creano le condizioni per tale strategia.

\bookdiagramplaceholder{assets/diagrams/image42.png}{42}

\emph{Difesa Nimzo Indiana:} \textbf{1.d4 ¤f6 2.c4 e6 3.¤c3 ¥b4 4.£c2 d5
5.¤xd5 exd5}

\emph{Difesa Gruenfeld:} \textbf{1.d4 ¤f6 2.c4 g6 3.¤c3 d5 4.¤f3 ¥g7
5.¥g5 ¤e4 6.¤xd5 ¤xg5 7.¤xg5 e6 8.¤f3 exd5}

\bookdiagramplaceholder{assets/diagrams/image43.png}{43}

\emph{Difesa Caro Kann:} \textbf{1.e4 c6 2.d4 d5 3.exd5 ¤xd5 4.¥d3 ¤c6
5.c3}

\bookdiagramplaceholder{assets/diagrams/image44.png}{44}

\emph{Difesa Cambridge Springs:} \textbf{1.d4 d5 2.c4 e6 3.¤c3 ¤f6 4.¥g5
c6 5.e3 ¤bd7 6.¤f3 £a5 7.¤xd5 exd5}

\bookdiagramplaceholder{assets/diagrams/image45.png}{45}

Il caso tipico, quello in cui l'attacco di minoranza avviene quasi di
norma, si ha nella Variante di Cambio del Gambetto di Donna:

\textbf{1.d4 d5 2.c4 e6 3.¤c3 ¤f6 4.¥g5 ¥e7 5.¤xd5 exd5 6.¤f3 c6}

\bookdiagramplaceholder{assets/diagrams/image46.png}{46}

Data la semplicità, e linearità di esecuzione, l'attacco di minoranza
può essere molto pericoloso anche in mano a un giocatore non
espertissimo.

Le reazioni possibili sono, in sostanza, due:

\emph{a)} contrattaccare sul lato opposto a quello in cui si sviluppa
l'attacco di minoranza;

\emph{b)} contrastarlo in modo diretto sullo stesso lato.

\textbf{a) Il contrattacco a un attacco di minoranza}

Abbiamo visto che nell'attacco di minoranza si spingono a fondo i propri
pedoni dopo avere ovviamente arroccato sull'altro lato della scacchiera.
La reazione energica sull'ala opposta avrà quindi le caratteristiche di
un assalto all'arrocco.

La partita che segue ne è un esempio e riassume bene quanto finora
spiegato.

Averbach--Konstantinopolsky, Mosca 1966 \textbf{(Gambetto di Donna
Rifiutato)}

\textbf{1.d4 d5 2.c4} si potrebbe discutere se a rigor di termini questo
debba essere considerato un gambetto. Dopo 2... £xc4 il vantaggio di
materiale è solo temporaneo non avendo modo il Nero di difendere il Pc4
a oltranza. Per esempio: 3.e3 b5 4.a4 c6 5. axb5 ¤xb5 6.£f3 con attacco
imparabile alla Torre in a8.

\textbf{2... c6 3.¤f3 e6 4.¤xd5 exd5 5.¤c3 ¤f6 6.¥g5 ¥e7 7.£c2 ¤bd7 8.e3
0-0 9.¥d3 ¦e8 10.0-0 ¤f8 11.¦ab1}

\bookdiagramplaceholder{assets/diagrams/image47.png}{47}

Finita la fase di sviluppo il Bianco prepara la spinta in b4. Ha inizio
l'attacco di minoranza. A mio parere il Bianco poteva arrivare alla
spinta in b4 un tempo prima giocando 11.¥xf6! (distoglie l'¥e7 dal
controllo di b4) con la seguente continuazione: 11... ¥xf6 12. b4

\textbf{11... ¤e4 12.¥xe7 £xe7 13. b4 a6 14.a4 ¤g6} alle spinte dei
pedoni sul lato di Donna il Nero reagisce imbastendo un contrattacco
sull'ala di Re.

\textbf{15.b5 axb5 16.axb5 ¥g4 17.¥xe4!}

\bookdiagramplaceholder{assets/diagrams/image48.png}{48}

In una partita giocata a Kiev nel 1954, in occasione del XXI Campionato
sovietico, Taimanov, contro Nezmetdinov, giocò meno correttamente 17.
¤d2? ma dopo 17... ¤xd2, 18.£xd2 ch4!, (minaccia 19... ¤f3+ o anche
19... ¥h3) l'attacco del Nero arriva prima. Per esempio: 19.¥e2 ¥h3!,
20.gxh3 £g5+, 21.¥g4 ¤f3+. Taimanov fu allora costretto a ripiegare su
19.f3 £xe3+ 20.£xe3 ¦xe3, 21.fxg4 ¦xd3, 22.bxc6 bxc6 (non si può 22...
¦xc3? per 23.¤xb7 ¦b8, 24.¦bc1 ¦c4, 25.¦xc4 £xc4, 26.¦a1), 23.¤e2 ¦d2,
24.¦f2 h6, 25.¦bf1 ¤g6, e riuscì a salvare questa partita solo grazie a
successivi errori del Nero.

\textbf{17... £xe4 18.¤d2! ¥f5} una mossa dubbia. Forse è migliore 18...
f5.

\textbf{19.bxc6 bxc6} lo scopo dell'attacco di minoranza è raggiunto: il
c6 pedone è isolato.

\textbf{20.¤e2 ch4} il Bianco rafforza il lato di Re e attacca il Pc6.

\textbf{21.¤g3 ¥g6}

\bookdiagramplaceholder{assets/diagrams/image49.png}{49}

Si imponeva 21... ¦ac8. Konstantinopolsky temeva 22. f3 senza accorgersi
della forte risposta 22... £g5. Ora il Pc6 cade senza sufficiente
compenso.

\textbf{22.£xc6 ¦ac8 23.£d5 f5 24.¦fc1 ¦xc1+ 25.¦xc1 f4 26.exf4 e3
27.£d3+ ¥f7 28.£xe3 £d7 29.C£e4 ¥g6 30.f5 ¤xf5 31.¤xf5 ¦xe4 32.¤d6 ¦xe3
33.fxe3 £d7 34.¦c8+ £xc8 35.¤xc8} il Nero abbandona.

\bookdiagramplaceholder{assets/diagrams/image50.png}{50}

La strategia del Bianco ha trionfato!

\textbf{b) contrastare l'attacco sullo stesso lato.}

Il modo più drastico di rispondere a un attacco di minoranza è, dopo
b2-b4, di giocare b7-b5 e piazzare un Cavallo in c4. La casa debole in
c5 trova il suo corrispettivo in quella c4 mentre il pedone c6 debole
non è facilmente sfruttabile dal Bianco dopo che il Nero, piazzato un
Cavallo in c4, ha chiuso la colonna a possibili attacchi di pezzi
pesanti. Si capisce pertanto come la mossa b7-b5 sia strettamente
connessa con ¤c4; se quest'ultima non fosse possibile la spinta del
pedone in b5 si rivelerebbe un gravissimo indebolimento. Attenzione
però: questa strategia può essere adottata solo se al Bianco è
interdetta la mossa b2-b3 che scaccerebbe il ¤c4 e renderebbe debole il
pedone in c6.

Mambrini--Leoncini, Imperia 1992 \textbf{(Gambetto di Donna Rifiutato)}

\textbf{1.d4 d5 2.c4 e6 3.¤c3 ¤f6 4.¥g5 ¥e7 5.¤xd5} il Bianco svela i
suoi piani. Per confermare la norma di Maestro era importante che
vincessi questa partita. Dover entrare in una variante in cui la
strategia del Bianco è semplice e valida per molte mosse, non mi fece
per niente piacere.

\textbf{5... exd5 6.£c2 g6!} nella partita di Donna lo sviluppo
dell'Alfiere campochiaro è un problema non irrilevante per il Nero. Con
la minaccia ¥f5 il Nero pone le basi per cambiarlo con quello buono del
Bianco. Il Bianco può evitare questa variante con l'adozione una diversa
successione di mosse; per esempio giocando 6.e3, in modo da replicare
alla spinta in g6 con 7.¥d3, mossa che impedisce 7... ¥f5, cui
seguirebbe 8.¥xf5! gxf5 e il Nero dove arrocca?

\textbf{7.e3 ¥f5 8.¥d3 ¥xd3} se il Bianco vuole evitare il cambio degli
Alfieri su casa chiara può giocare 8.£d3.

\textbf{9.£xd3 ¤bd7} una mossa importante. In caso di cambio in f6
occorre riprendere di Cavallo per non distogliere l'¥e7 dal controllo di
b4 (cfr. nota alla undicesima mossa bianca nella partita precedente).

\textbf{10.¤f3 0-0 11.0-0 c6 12. ¦ab1}

\bookdiagramplaceholder{assets/diagrams/image51.png}{51}

Eccoci
giunti al bivio: attaccare sul lato di Re o contrastare palmo a palmo
l'espansione del Bianco sull'ala di Donna? In questa partita optai per
la seconda strategia per ragioni puramente psicologiche. Avendo più
punti Elo del mio avversario volevo evitare che il suo compito fosse
facilitato dal dover eseguire un piano semplice.

Opporsi all'espansione significava costringerlo a una battaglia dalla
strategia non scontata dove speravo di far valere la mia presunta
(quanto è attendibile il punteggio Elo?) maggiore forza.

\textbf{12... a5 13.a3?} il Bianco deve arrivare alla spinta in b4 ma
dopo questa mossa tale spinta non è più possibile! Si dovevano muovere i
pedoni in un o¢dine diverso: b3, a3 e b4. La spinta in a5 è utile a
ritardare l'espansione bianca.

\textbf{13... a4!} ora a una eventuale spinta del pedone b, il Nero
cambia il pedone e il Bianco rimane con un pedone arretrato in a3.

\textbf{14.£c2 b5} solo ora che il Bianco non ha più modo di gua¢dare
con un pedone la casa c4 questa mossa può essere giocata.

\textbf{15.¥xf6 ¥xf6} il Nero può riprendere tranquillamente d'Alfiere
non essendo più necessario controllare la casa b4. Inoltre il Cavallo in
d7 è destinato ad andare in c4 (via b6).

\textbf{16.¤a2 ¦c8 17.¤b4} siccome in c6 c'è un pedone debole, il Bianco
ripiazza i pezzi per attaccarlo. Ma, come si vedrà, l'idea è
irrealizzabile. La casa ideale del Cavallo bianco è piuttosto in c5.

\textbf{17... ¤b6 18.b3?}

\bookdiagramplaceholder{assets/diagrams/image52.png}{52}

La mossa perdente! Il Bianco si accorge di non poter prendere in c6:
18.¤xc6? £d7!, 19. ¦fc1 ¤c4 e, dopo la ritirata del Cavallo bianco,
20... ¤xe3. Inoltre l'attacco di minoranza è stato frustrato sul
nascere, di più, è stato il Nero a guadagnare terreno sul lato di Donna.
In simili occasioni è facile sbagliare.

\textbf{18... axb3} il Nero ha finalmente un obiettivo concreto: il Pa3!

\textbf{19.£xb3 £d6 20.¦fc1 ¦c7} il Nero intende raddoppiare le Torri
sulla colonna ``a'' ma la ¦c8 per il momento non può essere distolta
dalla difesa del \textbf{§}c6, quindi lascia libero il passaggio verso
la colonna ``a'' alla ¦f8 e prepara il raddoppio in a7.

\textbf{21.¤a6 ¦a7 22.£d4 ¦d8 23.¤c5 T£a8 24.¤e4 £e6 25.¤c5} il Cavallo
si piazza nel buco in c5 ma, tutto sommato, a questo punto era meglio
cambiare in f6. Il mio avversario sperava forse che acconsentissi alla
patta per ripetizione di posizione.

\textbf{25... £e8! 26.¦a1 ¥e7 27.¤e5?} una mossa ottimista. Il Bianco
non sembra essersi reso conto che è giunto il momento di difendersi con
i denti. Sarebbe stato meglio 27.¤d2 per essere pronti a cambiare il
Cavallo nero appena spicca il salto in c4.

\textbf{27... f6 28. ¤ed3 ¤c4}

\bookdiagramplaceholder{assets/diagrams/image53.png}{53}

Finalmente il Cavallo si installa nella casa c4 e con effetti
dirompenti. Il terzo attacco al Pa3 è senza rimedio.

\textbf{29.¦xc4?} vistosi perduto il Bianco sacrifica la qualità
sperando in una boccata d'ossigeno. In realtà questa mossa affretta la
fine.

\textbf{29... bxc4 30.¤f4 ¦b8 31.£c3 ¥xc5 32.£xc5 £e5} in vantaggio di
materiale il Nero può cambiare tranquillamente i pezzi e, in caso di
ritirata della Donna bianca, questa mossa guadagna comunque spazio.

\textbf{33.£c1 ¦ab7} dopo questa mossa il Bianco non ha più difese. La
conclusione è rapida.

\textbf{34.f3 £xa1! 35.£xa1 ¦b1+ 36.£xb1 ¦xb1+ 37.¢f2 c3}

\bookdiagramplaceholder{assets/diagrams/image54.png}{54}

Il Pedone ``c'' marcia trionfalmente a promozione perciò il Bianco
abbandona.

\textbf{4. I pedoni forti}

\textbf{a) Il pedone passato isolato}

Per pedone passato (o libero) si intende un pedone che, davanti a sé, ha
sgombre di pedoni avversari la colonna in cui si trova e quelle
adiacenti.

La forza di un pedone libero ma isolato dipende da alcuni fattori:

\emph{a)} La sua forza aumenta più è avanzato. Se si trova nella seconda
o terza traversa, tra l'importanza strategica di essere libero e la
debolezza di essere isolato prevalgono in genere i fattori di debolezza.
L'avversario infatti ha di norma tutto il tempo e lo spazio per
preparare una casa di blocco ed attaccarlo in forze.

Se invece si trova ben difeso oltre la quarta traversa allora prevalgono
i motivi che ne fanno una forza. Oltre a rappresentare una minaccia
costante di promozione, può essere spinto, ed eventualmente sacrificato,
per sgomberare una casa o aprire una colonna. È vero che un pedone
libero ma isolato lega anche i pezzi alla sua difesa ma il vantaggio di
spazio e la maggiore libertà di manovra sono in genere un buon compenso.

\emph{b)} Il valore della sua forza aumenta in rapporto alla sua
mobilità. Per questa ragione il difendente deve bloccarne l'avanzata. Il
bloccatore più adatto sono i pezzi minori, perché non scacciabili con un
semplice attacco.

\emph{c)} La sua forza aumenta a mano a mano che ci si avvicina al
finale di partita. In finale la sua forza è massima quando è distante
dal resto dei pedoni.

\bookdiagramplaceholder{assets/diagrams/image55.png}{55}

Entrambi i giocatori hanno un pedone libero ma quello del Bianco vale di
più perché è lontano.

La strategia per vincere è semplice. Il Bianco spingerà il pedone ``a''
ed il Re Nero sarà costretto ad andare a prenderlo; in questo modo però
si allontanerà dai propri pedoni che saranno catturati facilmente dal
Bianco.

\textbf{1.a4 ¢c4} distratto dall'azione diversiva il cane da guardia
lascia sole le pecorelle che divengono facili prede del lupo bianco.

\textbf{2.a5 ¢c5 3.¢d3 ¢b5 4.¢xd4 ¢xa5 5.¢e5} il Bianco vince.

\textbf{b) Il pedone passato sostenuto}

Un pedone libero, e sostenuto da altri pedoni, è un obiettivo
importante. La sua forza d'attacco è ancora maggiore di quello isolato
perché non abbisogna di pezzi alla sua difesa, pezzi che così possono
essere usati in altri compiti, per essere sostenuto. Anche in questo
caso la sua forza aumenta nel finale ed è massima in quello di soli
pedoni perché limita drasticamente l'azione del Re avversario inchiodato
al suo controllo.

\bookdiagramplaceholder{assets/diagrams/image56.png}{56}

Il Bianco vince facilmente andando a prendere con il Re i pedoni
avversari.

Il Re Nero, non potendo uscire dal quadrato a5-a8-d5-d8, non ha alcun
modo di contrastare questa semplice manovra.

Abbiamo già avuto modo di vedere in azione la terribile forza di un
pedone passato (cfr. partita Sokolsky--Botvinnik dopo la
24\textsuperscript{a} mossa del Nero e anche la Mambrini--Leoncini dopo
la 29\textsuperscript{a} mossa).

Vediamo ancora qualche esempio.ù

\bookdiagramplaceholder{assets/diagrams/image57.png}{57}

Questa posizione si verificò nella partita giocata nel torneo A.V.R.O.
del 1938 tra Botvinnik e Capablanca.

Il Nero ha un pedone in più e la spinta del pedone libero ``a'' è un
piano semplice e vincente; il compenso del Bianco si basa sul
\textbf{§}e6 passato e avanzato. Il Nero lo ha bloccato con la Donna e
si appresta a far valere il pedone in più; purtroppo per lui la Donna
non è un bloccatore ideale perché al minimo attacco è costretta a
spostarsi. In apparenza comunque non si vede quale pezzo possa darle
fastidio.

\textbf{30.¥a3!! £xa3} la mossa del bianco è un fulmine a ciel sereno!

\textbf{31.ch5+!} scardina la posizione del Nero e confida nella forza
congiunta di Donna e pedone passato.

\textbf{31... gxh5 32.£g5+ ¢f8 33.£xf6+ ¢g8 34.e7} decisiva! Al Nero non
rimane che un inutile caccia allo scacco perpetuo.

\textbf{34... £c1+ 35.¢f2 £c2+ 36.¢g3 £d3+ 37.¢h4 £e4+ 38.¢xh5 £e2+
39.¢h4 £e4+ 40.g4 £e1+ 41.¢h5} non ci sono più scacchi pertanto il Nero
abbandona.

Il diagramma 58 riproduce la posizione prodottasi dopo la
ventiquattresima mossa del Nero della partita Bernstein--Capablanca del
torneo di Hastings del 1922.

\bookdiagramplaceholder{assets/diagrams/image58.png}{58}

Il pedone passato in c3 ha concentrato su di sé tutti i pezzi della
scacchiera. La sua presenza limita alquanto le possibilità del Bianco,
costretto a muoversi in spazi ristretti.

Il Bianco ha realizzato un blocco in c2 ma il bloccatore è un pezzo
pesante e quindi sarà costretto a muoversi in caso di un attacco di
Cavallo. In effetti, un eventuale raddoppio delle Torri nere sulla
colonna ``c'' consentirebbe al Nero di spostare il Cavallo che potrebbe
ripiazzarsi più convenientemente in b4 (o a3) per scacciare la Torre
dalla casa c2.

Il Bianco ha in animo di impedire il raddoppio e gioca:

\textbf{24.¤b3 ¦c6 25.¤d4} dove mettere la Torre? La casa c4 è
controllata dalla Donna; se la Torre torna in c5 il Bianco gioca il
Cavallo in b3 e ripete la posizione; a 25... ¦c7 segue invece 26. ¤b5 ed
il pedone c3 cade. Il Nero deve quindi andare in c8 e rita¢dare di varie
mosse l'attacco alla casa c2. Nel frattempo il Bianco può prendere
provvedimenti. Questo almeno è quanto doveva pensare Bernstein ma lo
attendeva una sorpresa:

\textbf{25... ¦c7} Bernstein, non vedendo il sottile tranello basato
sulla debolezza della prima traversa, deve aver pensato a un errore del
Campione del Mondo e, come preventivato, decide di guadagnare il pedone.

\textbf{26.¤b5 ¦c5 27.¤xc3 ¤xc3 28.¦xc3 ¦xc3 29.¦xc3 £d2!!} ecco la
bellissima mossa che Capablanca aveva preparato. Dopo questa mossa non
ci sono più difese. A 30.£xb2 segue 30... ¦d1 matto; se invece 30.¦c2
£d1+ 31.£f1 £xc2; se, infine, 30.£e1 £xc3! 31.£xc3 ¦d1 e matto alla
seguente.

\textbf{30.} il Bianco abbandona.

\bookdiagramplaceholder{assets/diagrams/image59.png}{59}

Per finire seguiamo le evoluzioni dei pezzi in una partita in cui
prevale la tecnica difensiva. Lo studio attento di questa istruttiva
partita dice molto su come contrastare la forza di un pedone passato.

Leonhardt--Nimzowitsch, San Sebastian 1912 \textbf{(Difesa Philidor)}

\textbf{1.e4 e5 2.¤f3 d6 3.d4 ¤f6 4.¤c3 exd4} forse è migliore non
cedere il centro e giocare 4... ¤bd7 5.¥c4 ¥e7 6.0-0 0-0 7.a4 c6.
Nimzowitsch fu un grande sperimentatore di linee difensive portando
molte innovazioni nel campo della teoria delle aperture e rivisitando in
chiave moderna, come fa in questa partita, difese considerate inferiori.

\textbf{5.¤xd4} una linea più aggressiva si ha con 5.£xd4 e se 5... ¤c6
allora 6.¥b5! ¥d7 7.¥xc6 e la coppia degli Alfieri trova compenso dal
vantaggio del Bianco al centro.

\textbf{5... ¥e7 6.¥e2 0-0 7.0-0 ¤c6 8.¤xc6 bxc6} Nimzowitsch commenta:
\emph{``Il cambio crea vantaggi per entrambi: il Nero ottiene una massa
di pedoni centrali più compatta (che impedisce l'occupazione della casa
d'avamposto d5 con l'eventuale ¤d5), ma il §a7 è isolato e anche c5 può
diventare debole, come nel testo.''}

\textbf{9.b3 d5 10.e5 ¤e8 11.f4 f5} la minaccia f4-f5 costringe il Nero
a questa spinta di pedone che dona al Bianco un pedone passato e
sostenuto in e5.

\textbf{12.¥e3 g6}

\bookdiagramplaceholder{assets/diagrams/image60.png}{60}

Come sappiamo il Nero deve bloccare il pedone passato nella casa e6 e il
modo migliore per farlo è con un pezzo leggero. I pezzi leggeri che
possono andare in e6 sono solo due l'¥c8 e il ¤e8, ma è facile vedere
che un ¥e6 svolgerebbe un ruolo del tutto passivo quindi è ovvio che
Nimzowitsch intenda portarci il Cavallo. Oltre a bloccare il pedone
passato, un Cavallo in e6 prepara le importanti spinte di rottura in d4
e g5.

Questa mossa si rende dunque necessaria per preparare la strada al
Cavallo che giungerà in e6 via g7.

\textbf{13.¤a4 ¤g7 14.£d2 £d7 15.£a5} oltre a fare pressione sul lato di
Donna questa mossa vuole impedire la manovra nera a5-¥a6 col conseguente
cambio degli Alfieri su casa chiara.

\textbf{15... ¤e6 16.¦ad1 ¦d8 17.¤c5?} una mossa ``naturale'' ma in
realtà un errore strategico. L'idea di occupare la casa c5 è giusta ma
il Bianco doveva portarci l'Alfiere, e solo dopo aver cambiato l'Alfiere
camposcuro, piazzarci il Cavallo. È importante rimuovere il bloccatore
in e6 e solo un Cavallo poteva riuscirci.

\textbf{17... ¥xc5} eliminato il Cavallo c5 il Bianco non ha più modo di
scacciare il bloccatore e di far valere, nel mediogioco, il pedone
passato. Ora l'iniziativa passa nelle mani del Nero.

\textbf{18.¥xc5 ¥b7 19.¦f3 ¢f7 20.¦h3 ¢g7 21.¦f1 ¦e8 22.¦hf3 ¦ad8 23.¦d1
a6} 23... £xa7? perde la Donna dopo 23... ¦a8 24.£xb7 ¦eb8.

\textbf{24.b4 ¢h8 25.£a3 ¦g8 26.£c3 ¦g7 27.¢h1 T£g8 28.¥e3 c5 29.¦g3}
secondo Schlechter era migliore 29.bxc5, col probabile seguito 29... d4
30.¦xd4! ¤xd4 31.¥xd4 ¥xf3 32.¥xf3 con due Alfieri e due pedoni per le
Torri.

\textbf{29... d4 30.£a3 g5 31.¥c4 gxf4 32.¥xe6}

\bookdiagramplaceholder{assets/diagrams/image61.png}{61}

Il Bianco elimina finalmente il bloccatore ma ormai è ta¢di per
capovolgere le sorti della partita. Oltretutto il Nero ha in se¢bo una
piccola sorpresa tattica.

\textbf{32... ¥xg2+! 33.¢g1} l'Alfiere è intoccabile. Se 33.¦xg2 £c6! e
se 33.¢xg2 £c6+ 34.¢f1 fxg3 35.¥xg8 gxh2 e vince. Ma anche così le case
bianche divengono molto deboli e la sorte del Bianco è segnata.

\textbf{33... £xe6 34.¥xf4 ¥b7 35.bxc5 £d5 36.c6 ¥xc6 37.¢f2 ¦xg3
38.hxg3 £g2+ 39.¢e1 ¥f3 40.£xa6 £g1+ 41.} il Bianco abbandona.

\textbf{5. L'attacco di maggioranza}

I cambi possono, con facilità, condurre ad asimmetrie nella
distribuzione dei pedoni sui due lati della scacchiera, e in modo tale
da creare una maggioranza di pedoni di un colore su di un lato e una
maggioranza di pedoni del colore opposto sull'altro lato.

\bookdiagramplaceholder{assets/diagrams/image62.png}{62}

Il diagramma mostra un caso tipico. Il Bianco dispone di tre pedoni
contro due sul lato di Donna e il Nero di quattro contro tre su quello
di Re. Nelle posizioni con arrocchi omogenei la teoria considera
vantaggioso avere la maggioranza sul lato opposto a quello in cui sono
posti i Re: in una posizione come quella del diagramma entrambi i
giocatori possono crearsi un pedone libero (il Bianco sulla colonna
``c'' e il Nero su quella ``e'') e, in tali casi, ha maggior valore
quello più distante dai Re. La sua cattura costringe, infatti, il Re ad
allontanarsi dai pedoni restanti che divengono così preda
dell'avversario.

Vediamo qualche esempio di apertura che, in pochi tratti, produce
maggioranze su lati opposti.

\emph{Siciliana:} \textbf{1.e4 c5 2.c3 d5 3.exd5 £xd5 4.d4 ¤f6 5.¤f3 e6
6.¥d3 ¥e7 7.£xc5}

\bookdiagramplaceholder{assets/diagrams/image63.png}{63}

\emph{Caro Kann:} \textbf{1.e4 c6 2.d4 d5 3.exd5 ¤xd5 4.c4 ¤f6 5.¤c3 e6
6.c5}

\bookdiagramplaceholder{assets/diagrams/image64.png}{64}

\emph{Difesa Tarrasch:} \textbf{1.d4 d5 2.c4 e6 3.¤c3 c5 4.¤xd5 exd5
5.¤f3 ¤c6 6.g3 c4}

\bookdiagramplaceholder{assets/diagrams/image65.png}{65}

Non si creda che, raggiunta una maggioranza sul lato opposto a quello
dei Re, si sia già vinto da partita. Certo, il finale di soli pedoni è
favorevole ma al finale bisogna arrivarci! La maggioranza può anche
essere usata in centro partita per imbastire un attacco finalizzato alla
formazione di un pedone libero come succede in questa breve partita in
cui la strategia del Bianco trionfa.

Hort--Ciric, Amsterdam 1970 \textbf{(Difesa Caro Kann)}

\textbf{1.e4 c6 2.d4 d5 3.exd5 ¤xd5 4.c4} con l'attacco Panov il Bianco
evita continuazioni squisitamente posizionali tipiche della Caro Kann.
L'alternativa è 4.¥d3 col possibile seguito 4... ¤c6 5.c3 ¤f6 6.¥f4 ¥g4
7.£d3 £c8 8.¤d2 e6 9.¤gf3 ¥e7 10.0-0 0-0 11.¤e5 ¥h5 12.¦fe1 ¤xe5 13.¥xe5
£c6 con gioco equilibrato (Szapik-Seirawan, Olimpiadi di Malta 1980).

\textbf{4... ¤f6 5.¤c3 e6 6.¤f3 ¥e7 7.¥g5} si può anche giocare subito
7.c5 0-0 8.¥d3 b6 9.b4 a5 10.¤a4 con complicazioni.

\textbf{7... 0-0 8.c5} il Bianco si crea una maggioranza sul lato di
Donna e minaccia di spingere a fondo i propri pedoni. Al Nero è
richiesta una difesa precisa.

\textbf{8... b6} penso sia meglio giocare 8... ¤c6 come nella partita
Liberzon-Zaitsev del campionato sovietico del 1969, cui seguì 9.¦c1 ¤e4
10.¥xe7 £xe7 11.¥e2 ¦d8 12.0-0 e5 13.¤xe5 ¤xe5 14.£xe5 ¤xc3 15.¦xc3 d4
con gioco equivalente.

\textbf{9.b4 a5} è ovvio che, avendo deciso di spingere il pedone ``c''
in c5 il Bianco si è anche impegnato a sostenerlo.

\textbf{10.a3 ¤e4 11.¥xe7 £xe7 12.¤xe4!} è sbagliata 12.£c1? per 12...
axb4 13.¦xa1 £xa1 14.£xa1 ¤c6 15.£d2 bxc5 16.bxc5 £d7 con vantaggio del
Nero.

\textbf{12... £xe4 13.¤e5 ¦d8?} è meglio 13... ¥b7

\textbf{14.¥b5! axb4 15.axb4 ¦xa1 16.£xa1 bxc5 17.£xc5}

\bookdiagramplaceholder{assets/diagrams/image66.png}{66}

Con la comparsa di due pedoni liberi sul lato di Donna, il Bianco ha
raggiunto l'obiettivo strategico che si era proposto. Ora si tratta di
mobilitarli.

\textbf{17... £d7 18.£a5} una mossa precisa, l'attacco alla ¦d8 difende
in modo indiretto il ¤e5. Per esempio, se il Bianco avesse giocato
18.£a4? il Nero avrebbe reagito attaccando il ¤e5 con 18... £d5! e su
18.¤c4 avrebbe avuto il tempo di giocare 19... e3!

\textbf{18... £d5 19.0-0 ¦f8 20.£a1 ¦d8 21.¥c4 £d4 22.c6 £d6 23.b5} il
Bianco ha spostato in avanti la catena dei pedoni; la vittoria è vicina.

\textbf{23... £c7 24.£d2 £d6 25.£c2 f5 26.c7} il Nero abbandona.

\bookdiagramplaceholder{assets/diagrams/image67.png}{67}

A 26... £xc7 segue 27.¥xe6+!
